\textbf{Phần 1. Trắc nghiệm lựa chọn}
\setcounterpageref{ex}{0}

\begin{ex} %[4D3V12]
Lớp triệu gồm các hàng nào?
\choice
{\True Hàng triệu, hàng chục triệu, hàng trăm triệu}
{Hàng nghìn, hàng chục nghìn, hàng trăm nghìn}
{Hàng trăm, hàng nghìn, hàng triệu}
{Hàng triệu, hàng tỉ, hàng nghìn tỉ}
\loigiai{
Lớp triệu gồm: hàng triệu, hàng chục triệu, hàng trăm triệu.\\
\textbf{Đáp án của câu hỏi là A}
}
\end{ex}

\begin{ex} %[4D3V12]
Số \textit{Ba mươi lăm triệu} được viết là
\choice
{$3\,500\,000$}
{$350\,000\,000$}
{\True $35\,000\,000$}
{$35\,000$}
\loigiai{
Số đó là $35\,000\,000$.\\
\textbf{Đáp án của câu hỏi là C}
}
\end{ex}

\begin{ex} %[4D3V12]
Trong số $45\,678\,901$, chữ số $4$ thuộc hàng nào?
\choice
{Hàng triệu}
{\True Hàng chục triệu}
{Hàng trăm triệu}
{Hàng trăm nghìn}
\loigiai{
Chữ số 4 ở hàng chục triệu.\\
\textbf{Đáp án của câu hỏi là B}
}
\end{ex}

\begin{ex} %[4D3V12]
Số nào sau đây có 8 chữ số?
\choice
{$5\,432\,100$}
{$123\,456\,789$}
{$999\,999$}
{\True $12\,345\,678$}
\loigiai{
$12\,345\,678$ có 8 chữ số.\\
\textbf{Đáp án của câu hỏi là D}
}
\end{ex}

\begin{ex} %[4D3V12]
Số gồm $5$ trăm triệu, $2$ chục triệu và $8$ triệu viết là
\choice
{\True $528\,000\,000$}
{$52\,800\,000$}
{$502\,800\,000$}
{$520\,800\,000$}
\loigiai{
Số đó là $528\,000\,000$.\\
\textbf{Đáp án của câu hỏi là A}
}
\end{ex}

\begin{ex} %[4D3V12]
Trong số $123\,456\,789$, các chữ số thuộc lớp triệu là
\choice
{$4; 5; 6$}
{$7; 8; 9$}
{\True $1; 2; 3$}
{$1; 4; 7$}
\loigiai{
Lớp triệu gồm: 1, 2, 3.\\
\textbf{Đáp án của câu hỏi là C}
}
\end{ex}

\begin{ex} %[4D3V12]
Dân số của một quốc gia là $98\,500\,000$ người. Số này đọc là
\choice
{Chín mươi tám nghìn năm trăm người}
{\True Chín mươi tám triệu năm trăm nghìn người}
{Chín trăm tám mươi lăm triệu người}
{Chín mươi tám triệu năm mươi nghìn người}
\loigiai{
Đọc là: Chín mươi tám triệu năm trăm nghìn người.\\
\textbf{Đáp án của câu hỏi là B}
}
\end{ex}

\begin{ex} %[4D3V12]
Số liền sau của số $99\,999\,999$ là
\choice
{$99\,999\,998$}
{$100\,000\,001$}
{$10\,000\,000$}
{\True $100\,000\,000$}
\loigiai{
$99\,999\,999 + 1 = 100\,000\,000$.\\
\textbf{Đáp án của câu hỏi là D}
}
\end{ex}

\begin{ex} %[4D3V12]
Giá trị của chữ số 7 trong số $742\,500\,123$ là
\choice
{\True $700\,000\,000$}
{$7\,000\,000$}
{$70\,000\,000$}
{$700\,000$}
\loigiai{
Chữ số 7 ở hàng trăm triệu nên có giá trị là $700\,000\,000$.\\
\textbf{Đáp án của câu hỏi là A}
}
\end{ex}

\begin{ex} %[4D3V12]
Số gồm $4$ chục triệu, $5$ trăm nghìn và $6$ trăm viết là
\choice
{$45\,000\,600$}
{$40\,050\,600$}
{\True $40\,500\,600$}
{$450\,000\,600$}
\loigiai{
$40\,000\,000 + 500\,000 + 600 = 40\,500\,600$.\\
\textbf{Đáp án của câu hỏi là C}
}
\end{ex}

\begin{ex} %[4D3V12]
Một công ty đạt doanh thu $120$ triệu đồng trong quý 1, $150$ triệu đồng trong quý 2. Tổng doanh thu 2 quý viết bằng số là
\choice
{$27\,000\,000$ đồng}
{\True $270\,000\,000$ đồng}
{$2\,700\,000\,000$ đồng}
{$270\,000$ đồng}
\loigiai{
$120\,000\,000 + 150\,000\,000 = 270\,000\,000$ đồng.\\
\textbf{Đáp án của câu hỏi là B}
}
\end{ex}

\begin{ex} %[4D3V12]
Số nhỏ nhất có 9 chữ số khác nhau là
\choice
{$100\,000\,000$}
{$123\,456\,789$}
{$102\,345\,687$}
{\True $102\,345\,678$}
\loigiai{
Số nhỏ nhất có 9 chữ số khác nhau là $102\,345\,678$.\\
\textbf{Đáp án của câu hỏi là D}
}
\end{ex}

\textbf{Phần 2. Trắc nghiệm đúng sai}
\setcounterpageref{ex}{0}

\begin{ex} %[4D3V12]
Xét tính đúng sai về các số trong phạm vi lớp triệu:
\choiceTF
{\True $10\,000\,000$ đọc là mười triệu.}
{\True $100\,000\,000$ đọc là một trăm triệu.}
{Số $45\,000\,000$ có 7 chữ số.}
{\True Số $45\,000\,000$ có 8 chữ số.}
\loigiai{
\begin{itemize}
    \item $10\,000\,000$: mười triệu. Đúng.
    \item $100\,000\,000$: một trăm triệu. Đúng.
    \item $45\,000\,000$ có 8 chữ số (không phải 7 chữ số). Ý 3 Sai, Ý 4 Đúng.
\end{itemize}
\begin{itemize}
    \item a) \textbf{Đúng}.
    \item b) \textbf{Đúng}.
    \item c) \textbf{Sai}.
    \item d) \textbf{Đúng}.
\end{itemize}
}
\end{ex}

\begin{ex} %[4D3V12]
Cho số $250\,360\,000$. Xét tính đúng sai:
\choiceTF
{\True Lớp triệu của số đó là $250$.}
{\True Lớp nghìn của số đó là $360$.}
{Lớp đơn vị của số đó là $100$.}
{\True Số đó gồm $2$ trăm triệu, $5$ chục triệu, $3$ trăm nghìn và $6$ chục nghìn.}
\loigiai{
\begin{itemize}
    \item Lớp triệu: 250. Đúng.
    \item Lớp nghìn: 360. Đúng.
    \item Lớp đơn vị là 000 (không phải 100). Sai.
    \item Cấu tạo số chính xác. Đúng.
\end{itemize}
\begin{itemize}
    \item a) \textbf{Đúng}.
    \item b) \textbf{Đúng}.
    \item c) \textbf{Sai}.
    \item d) \textbf{Đúng}.
\end{itemize}
}
\end{ex}

\textbf{Phần 3. Trả lời ngắn}
\setcounterpageref{ex}{0}

\begin{ex} %[4D3V12]
Số bé nhất có 7 chữ số là số nào?
\shortans{1000000}
\loigiai{
\textbf{Đáp số:} $1000000$\\
Số đó là $1\,000\,000$.
}
\end{ex}

\begin{ex} %[4D3V12]
Viết số: Tám mươi triệu bốn trăm nghìn.
\shortans{80400000}
\loigiai{
\textbf{Đáp số:} $80400000$\\
Số đó là $80\,400\,000$.
}
\end{ex}

\begin{ex} %[4D3V12]
Trong số $567\,890\,123$, tính tổng các chữ số thuộc lớp triệu.
\shortans{18}
\loigiai{
\textbf{Đáp số:} $18$\\
$5 + 6 + 7 = 18$.
}
\end{ex}

\begin{ex} %[4D3V12]
Số liền trước của số $50\,000\,000$ là số nào?
\shortans{49999999}
\loigiai{
\textbf{Đáp số:} $49999999$\\
$50\,000\,000 - 1 = 49\,999\,999$.
}
\end{ex}

\begin{ex} %[4D3V12]
Tính giá trị chữ số 8 trong số $845\,000\,000$ theo đơn vị triệu (chỉ ghi số triệu).
\shortans{800}
\loigiai{
\textbf{Đáp số:} $800$\\
$800$ triệu.
}
\end{ex}

\begin{ex} %[4D3V12]
Tìm số có 8 chữ số lớn nhất có tất cả các chữ số thuộc lớp triệu bằng 9, lớp nghìn bằng 8 và lớp đơn vị bằng 7.
\shortans{99888777}
\loigiai{
\textbf{Đáp số:} $99888777$\\
Số đó là $99\,888\,777$.
}
\end{ex}
